% from ega1-2-fr.tex
\oldpage[I]{103}
Soient $(X,\mathscr{O}_X)$ un schéma local dont l'anneau est un corps $K$,
$\xi$ l'unique point de $X$, $(Y,\mathscr{O}_Y)$ un préschéma. Tout morphisme
$u:(X,\mathscr{O}_X)\to(Y,\mathscr{O}_Y)$ se factorise de façon unique en
$X\to\operatorname{Spec}(k(\psi(\xi)))\to Y$, où la seconde flèche désigne le
morphisme canonique, et la première correspond à un monomorphisme
$k(\psi(\xi))\to K$. Cela établit une correspondance biunivoque canonique
entre l'ensemble des morphismes $(X,\mathscr{O}_X)\to(Y,\mathscr{O}_Y)$ et
l'ensemble des monomorphismes $k(y)\to K$ $(y\in Y)$.
\end{corollary}

\begin{corollary}[2.4.7]
\label{I.2.4.7-fr}
Pour tout $y\in Y$, tout morphisme canonique
$\operatorname{Spec}(\mathscr{O}_y/\mathfrak{a}_y)\to Y$ est un monomorphisme
d'espaces annelés.
\end{corollary}

On l'a déjà vu lorsque $\mathfrak{a}_y=0$
(\hyperref[I.2.4.2-fr]{2.4.2}), et il suffit d'appliquer
(\hyperref[I.1.7.5-fr]{1.7.5}).

\begin{remark}[2.4.8]
\label{I.2.4.8-fr}
Soient $X$ un schéma local, $a$ son unique point fermé. Comme tout ouvert
affine contenant $a$ est nécessairement $X$ tout entier, tout
$\mathscr{O}_X$-Module \emph{inversible}
(\hyperref[0.5.4.1-fr]{0, 5.4.1}) est nécessairement \emph{isomorphe à
$\mathscr{O}_X$} (ou, comme on dit encore, est \emph{trivial}). Cette propriété
n'a pas lieu en général pour un schéma affine quelconque
$\operatorname{Spec}(A)$~; on verra au chap.~V que si $A$ est un anneau normal,
elle est vraie lorsque $A$ est \emph{factoriel}.
\end{remark}

\subsection{Préschémas au-dessus d'un préschéma}
\label{subsection:I.2.5-fr}

\begin{definition}[2.5.1]
\label{I.2.5.1-fr}
Étant donné un préschéma $S$, on dit que la donnée d'un préschéma $X$ et d'un
morphisme de préschémas $\varphi:X\to S$ définit un préschéma $X$ \emph{au-dessus
du préschéma $S$}, ou un \emph{$S$-préschéma}~; on dit que $S$ est le
\emph{préschéma de base} du $S$-préschéma $X$. Le morphisme $\varphi$ est appelé
le \emph{morphisme structural} du $S$-préschéma $X$. Lorsque $S$ est un schéma
affine d'anneau $A$, on dit aussi que $X$ muni de $\varphi$ est un préschéma
au-dessus de l'anneau $A$ (ou un $A$-préschéma).
\end{definition}

Il résulte de (\hyperref[I.2.2.4-fr]{2.2.4}) que la donnée d'un préschéma
au-dessus d'un anneau $A$ équivaut à la donnée d'un préschéma
$(X,\mathscr{O}_X)$ dont le faisceau structural $\mathscr{O}_X$ est un faisceau
de $A$-\emph{algèbres}. \emph{Un préschéma quelconque peut donc toujours être
considéré de façon unique comme un $\mathbf{Z}$-préschéma.}

Si $\varphi:X\to S$ est le morphisme structural d'un $S$-préschéma $X$, on dit
qu'un point $x\in X$ est \emph{au-dessus d'un point $s\in S$} si
$\varphi(x)=s$. On dit que $X$ \emph{domine} $S$ si $\varphi$ est un morphisme
dominant (\hyperref[I.2.2.6-fr]{2.2.6}).

\begin{env}[2.5.2]
\label{I.2.5.2-fr}
Soient $X$, $Y$ deux $S$-préschémas~; on dit qu'un morphisme de préschémas
$u:X\to Y$ est un \emph{morphisme de préschémas au-dessus de $S$} (ou un
$S$-\emph{morphisme}) si le diagramme
\phantomsection\label{I.2.5.2.diagram-fr}
\[
  \xymatrix{
    X\ar[rr]^u\ar[dr] & & Y\ar[dl]\\
    & S &
  }
\]
(où les flèches obliques sont les morphismes structuraux) est commutatif~:
cela entraîne que pour tout $s\in S$ et tout $x\in X$ au-dessus de $s$, $u(x)$
doit aussi être au-dessus de $s$.
\end{env}

Cette définition montre aussitôt que le composé de deux $S$-morphismes est un
$S$-morphisme~; les $S$-préschémas forment donc une \emph{catégorie}.

On notera $\operatorname{Hom}_S(X,Y)$ l'ensemble des $S$-morphismes d'un
$S$-préschéma $X$ dans un $S$-préschéma $Y$~; le morphisme identique d'un
$S$-préschéma $X$ sera désigné par $1_X$.

Lorsque $S$ est un schéma affine d'anneau $A$, on dira aussi
$A$-\emph{morphisme} au lieu de $S$-morphisme.

\begin{env}[2.5.3]
\label{I.2.5.3-fr}
